% The abstract of the notes.  Three parts, in this order: Översikt (what
% the deck does and where it goes), Lärandemål (the learning objectives),
% Förkunskaper (what the student should have met before).  A Läsning part
% is optional.  Section headings are \emph, not \section: the abstract is
% one block of prose.

\emph{Översikt}
Den här övningen är till för att göra det själv.  Föreläsningarna
\emph{Variabler och utskrifter} och \emph{Funktioner} visade vad en
tilldelning gör, vad ett anrop gör och hur ett program som upprepar sig
går att dela upp; här är det du som skriver programmen.  Uppgifterna
börjar i ett program du bara ska läsa och förutsäga, går vidare till två
vardagssysslor du ska dela upp i namngivna delar --- en där delarna
\emph{gör} något och en där de \emph{lämnar tillbaka} något --- och
fortsätter till en beräkning där tre olika formler delar på ett och samma
sista steg.  Sist får du läsa ett program som någon annan har skrivit och
säga vad som gör det svårt att läsa, ändra och använda.  Varje uppgift
följs av ett lösningsförslag, och det är ett förslag: nästan varje uppgift
går att lösa på flera rimliga sätt.

\emph{Lärandemål}
Efter den här övningen ska studenten kunna:
% Lo03
\begin{restatable}{lo}{OvningVariablesLODefine}\label{OvningVariablesLODefine}%
  Skriva egna funktioner med \mintinline{python}{def}, med parametrar och
  returvärde, och anropa dem.
\end{restatable}
% Lo03
\begin{restatable}{lo}{OvningVariablesLOCall}\label{OvningVariablesLOCall}%
  Följa exekveringen genom ett program där funktioner anropar varandra,
  och säga i vilken ordning utskrifterna kommer.
\end{restatable}
% Lo01, Lo02, Lo03
\begin{restatable}{lo}{OvningVariablesLODecompose}%
  \label{OvningVariablesLODecompose}%
  Dela upp ett program i funktioner så att kod inte upprepas, och namnge
  de principer uppdelningen följer: DRY, SRP och KISS.
\end{restatable}
% Lo07
\begin{restatable}{lo}{OvningVariablesLODocstring}%
  \label{OvningVariablesLODocstring}%
  Dokumentera en funktion med en
  \foreignlanguage{english}{docstring} enligt PEP~257.
\end{restatable}
% Lo07
\begin{restatable}{lo}{OvningVariablesLONames}\label{OvningVariablesLONames}%
  Välja beskrivande namn på funktioner och variabler och markera
  konstanter enligt PEP~8.
\end{restatable}
% Lo05
\begin{restatable}{lo}{OvningVariablesLOOutput}\label{OvningVariablesLOOutput}%
  Skriva ut text och värden med \mintinline{python}{print} och
  f-strängar så att utskriften blir läsbar för den som kör
  programmet.
\end{restatable}
% Lo11
\begin{restatable}{lo}{OvningVariablesLOReview}\label{OvningVariablesLOReview}%
  Granska ett program som någon annan har skrivit och peka ut vad som gör
  det svårt att läsa, ändra och använda.
\end{restatable}

\emph{Förkunskaper}
Studenten bör ha tagit del av föreläsningen \emph{Variabler och
utskrifter}: hon ska kunna följa ett pythonprogram uppifrån och ned, sats
för sats, veta vad en tilldelning gör och kunna skriva ut värden med
\mintinline{python}{print} och f-strängar.  Studenten bör också ha tagit
del av föreläsningen \emph{Funktioner}: övningen förutsätter
\mintinline{python}{def}, parametrar, argument och
\mintinline{python}{return}, att en funktion kan anropa en annan,
och de tre principerna DRY, SRP och KISS.
Från föreläsningen \emph{Algoritmiskt tänkande} förutsätts stegvis
förfining --- att ge en grupp steg ett namn och förfina namnet för sig ---
och principen DRY.

\emph{Läsning}
Namngivnings- och formateringsreglerna står i PEP~8\autocite{pep8} och
konventionen för \foreignlanguage{english}{docstrings} i
PEP~257\autocite{pep257}.  Du behöver inte kunna dem utantill, men det är
värt att veta vad som står där.
